\subsection{The matching engine}
\label{sec:ha:disabled:me}

The matching engine dies and is restarted. No member's connection is broken and
the venue can still take orders; what stops is anything being done with them.

\subsubsection{What state it holds}
\label{sec:ha:disabled:me:state}

\textbf{The set of open orders.} Which orders the venue has accepted and not yet
cancelled or rejected, and the session each belongs to. This is what lets an
order identifier be recognised as one already in use, and a cancel be answered
rather than guessed at. It is the only copy: the sequencer holds the record of
what was taken and in what order, which is a different thing and cannot be
turned into this one cheaply.

\textbf{How far that set has been brought.} The sequence number up to which the
set reflects what the sequencer has committed. Without it the set is a
collection of orders with no way to say what is missing from it.

\subsubsection{Where the copy that outlives it lives}
\label{sec:ha:disabled:me:durable}

In a memory-mapped region of fixed-size records, one for each open order, held
beside the engine and separate from the structures it works with. An order is
written to it as it is accepted and taken out of it as it is cancelled, and it is
read back when the engine starts. See
\texttt{docs/durability/open\_order\_checkpoint.md}.

It survives the process dying and not the machine dying. That is the guarantee
this configuration is specified to give, and the enabled configuration's answer to
a machine is the other machine.

\textbf{Why a region and not the sequencer's log}, which looks as though it should
serve. The log records changes rather than state, so rebuilding the set from it
means replaying from the beginning --- and the earlier part of the log is
discarded as it grows.

\begin{gap}{BUG-0048}
  The history retained is whatever has not yet been discarded, which is a
  quantity of records rather than a period of time. Rebuilding a set of orders
  of any age from the log is therefore not possible, however the rebuilding is
  written.
\end{gap}

\textbf{What is required instead is a checkpoint}: the set itself, recording the
orders the component holds open and the sequence number they are current to.
That is what the region is, and it is the copy that outlives the
process, and the log supplies only what has happened since.

\subsubsection{How a restarted instance becomes current}
\label{sec:ha:disabled:me:recovery}

It reads the checkpoint, rebuilds the set as it stood at the sequence number the
checkpoint records, and asks the sequencer for everything committed after that.
It applies what comes back without issuing execution reports for it --- those
were issued the first time --- and begins accepting orders once it has caught
up.

\textbf{This is the same operation a promoted instance performs.} A follower
that is promoted holds a replica of the set and knows the sequence number it
reached; it asks for everything after that, applies it silently, and begins
accepting. The two differ only in where the starting number comes from --- a
file in one case, a replica in the other.

Two things follow. The configuration with high availability enabled needs no
separate promotion path, because promotion is this recovery starting from a
number learned another way. And because the same mechanism runs on every restart,
it is exercised continually rather than only when a test produces a failover.


\subsubsection{What the member is told}
\label{sec:ha:disabled:me:member}

\begin{requirement}{R-0028}{Every accepted order is reported to its member, once}
  An order the matching engine accepted shall have an execution report reach the
  member that placed it, whatever restarted between the acceptance and the
  delivery, and the member shall not be led to believe the acceptance happened
  twice.
  \rationale{Catching up after a restart applies orders without issuing reports,
  on the ground that the reports were issued the first time. An order accepted
  whose report never left the engine is exactly the case where that ground does
  not hold, so something must be able to tell the two apart.}
  \uncovered{no scenario restarts the engine between acceptance and delivery}
\end{requirement}

\begin{requirement}{R-0029}{A cancel request is always answered}
  A request to cancel shall receive either a report saying the order was
  cancelled or a rejection saying why, whatever restarted while the request was
  in flight.
  \rationale{A member that does not know whether its cancel took effect must
  assume the order is live and can hedge in neither direction. A rejection it
  can act on; silence it cannot.}
  \uncovered{no scenario restarts the engine with a cancel in flight}
\end{requirement}

\begin{requirement}{R-0030}{With no matching engine, the venue refuses rather than acknowledges}
  Where no matching engine exists to process orders, the venue shall refuse
  orders and tell members why, rather than accept orders it cannot process.
  \rationale{An order acknowledged and never processed is worse for a member
  than one refused: it believes the order live, may have hedged against it, and
  has nothing to cancel.}
  \uncovered{scenario 42 covers refusal, and not a restart}
\end{requirement}

\begin{requirement}{R-0018}{An open order survives a restart of the matching engine}
  An order the venue held open before the restart shall be held open afterwards,
  on the same terms, and shall remain cancellable by the member that placed it.
  \rationale{The member was told the order was accepted and has no reason to
  think a venue process has restarted. An order that quietly ceases to exist is
  the one failure a member cannot detect, cannot cancel and cannot hedge.}
  \covers{ha_test:50}
\end{requirement}

\begin{requirement}{R-0120}{A cancel for an order the venue no longer holds says which case it is}
  Where a member asks to cancel an order the venue does not currently hold, it
  shall be told whether the order once existed and is finished, or was never
  known to the venue at all.
  \rationale{These are different facts about the member's own position and it
  cannot establish either without being told. The protocol distinguishes them for
  that reason.

  There is a second reason particular to this venue. \emph{Unknown order} is also
  what a member is told when the venue has lost what it was holding --- so
  answering both cases the same way makes a venue that has forgotten a member's
  orders indistinguishable, from that member's side, from one answering
  correctly. The distinction is what stops the failure this chapter exists to
  prevent from hiding inside an ordinary reply.}
  \uncovered{no scenario cancels an order that has already been cancelled}
\end{requirement}

\begin{gap}{}
  A cancel for an order the venue does not currently hold is answered as an order
  it does not recognise, whether the order was cancelled a moment ago or never
  existed. A member cannot tell those apart, and cannot tell either of them from
  the venue having lost what it held.
\end{gap}

\begin{requirement}{R-0116}{The venue recovers what it held before deciding what to do with it}
  A restarted instance shall rebuild the set of orders it held before deciding
  whether to resume trading or to cancel them.
  \rationale{Cancelling requires knowing which orders to cancel and which member
  to tell about each. A venue that cancels without having recovered cannot report
  anything: it simply stops mentioning the orders, which is what \reqref{R-0020}
  exists to prevent. Recovery is therefore needed for both outcomes, and is not
  an alternative to cancelling but the thing that makes cancelling possible.}
  \covers{ha_test:51}
\end{requirement}

\begin{requirement}{R-0123}{A member is told what it had, even when the venue's own record of it is unusable}
  Where the record a component keeps of the orders it held cannot be used ---
  absent, damaged, or written by a different build --- the venue shall establish
  what was open from its record of what it took, cancel each order, and report
  the cancellation to the member that placed it. Where it cannot establish that
  either, it shall halt and say that it cannot account for what it was holding.
  \rationale{\reqref{R-0102} says such a record is not to be used. That leaves
  the venue unable to name what it held, and it cannot cancel what it cannot
  name --- so without a second route it would fall silent about every order,
  which is the outcome \reqref{R-0020} exists to prevent.

  The record of what the venue took is that second route. Rebuilding from it is
  far slower than reading a checkpoint, which is why it is not the ordinary
  recovery; but this is not the ordinary case, it happens once, and what the time
  buys is every member being told what became of its orders.

  Admitting that it cannot say is the last resort and is still better than
  silence: a member that is told the venue has lost track can act, where a member
  that is told nothing cannot.}
  \covers{ha_test:52}
\end{requirement}

\begin{requirement}{R-0117}{Orders open across a long absence are cancelled, reported, and trading halts}
  Where the venue was unable to match for longer than a stated period and orders
  were open when it stopped, it shall cancel each of them, send the member that
  placed it a report saying so, and halt. Where nothing was open, it shall resume
  without halting.
  \rationale{\textbf{The period is the length of the outage, not the age of the
  order.} An order that has rested since the morning in a working market is not a
  problem: its owner could have cancelled it at any moment and did not. What
  makes an order dangerous is that its owner was \emph{locked out of it} --- the
  venue could not match, so the member could not cancel, and the market moved
  while it could do nothing. The order is now priced for a market that has gone,
  and its owner was denied every opportunity to act on that.

  So the hazard has nothing to do with how long the order had been resting, and
  nothing to do with whether it was recoverable: a correctly recovered order that
  its member was locked out of for ten minutes is exactly as dangerous as one
  recovered from a corrupt region would have been.

  Cancelling removes the uncertainty, the reports tell each member where it
  stands, and the halt says the venue is not simply carrying on as though nothing
  had happened. The exception for an empty book matters because traffic is
  bursty: a long absence during a lull leaves nothing stale, nothing to report,
  and halting would achieve only an outage.

  Unlike a condition the venue cannot detect, both parts of this one are plain
  --- how long it was absent, and what it cancelled --- so it can be acted on
  without a person present. Lifting it still requires one, under
  \reqref{R-0023}.}
  \covers{ha_test:51}
\end{requirement}

\begin{requirement}{R-0118}{How long an absence may be before orders are cancelled is a stated figure}
  The period in \reqref{R-0117} shall be a figure the venue states, measured from
  the moment the venue stopped being able to match to the moment it resumes, and
  the venue shall determine that without depending on anything outside the
  component concerned.
  \rationale{It is the whole absence and not the recovery step. Noticing the
  death, restarting the process and rebuilding the set are all terms in it, and
  the restart alone was measured at about two and a half seconds --- so a figure
  configured against recovery time is set against the wrong quantity and will be
  reached later than intended.

  \textbf{It is measured from when matching stopped, not from the last order
  matched.} Traffic comes in bursts, so a quiet period before a failure is
  ordinary; taking the last order as the start would read a ten-minute lull
  followed by a two-second restart as a twelve-minute absence, and cancel every
  member's orders for having been quiet.

  \textbf{And the component must be able to work it out alone.} A supervisor
  knows exactly when it killed a process and when it started the replacement, and
  is the obvious source --- but nothing in this chapter may depend on a supervisor
  being there, and a component started by hand, in a debugger, hours after the
  last one stopped, must reach the same answer. It does: an hour has passed, the
  orders are an hour stale, and cancelling them is right.

  It is also a trading decision rather than a measurement: it says how long a
  member's resting order stays meaningful in this market, which is a judgement
  about members and not about how long the code takes.}
  \covers{ha_test:51}
\end{requirement}

\begin{requirement}{R-0036}{A cancelled order does not become open again}
  An order the member cancelled before a restart shall not be open after it.
  \rationale{Recovery reads a record of the state and then applies what followed
  it. A cancellation that falls between the two is the case that resurrects an
  order its member killed --- and the member, having been told it was cancelled,
  is not watching for it.}
  \uncovered{no scenario cancels an order and then restarts the engine}
\end{requirement}

\begin{requirement}{R-0020}{An order the venue cannot recover is cancelled to its member}
  Where the set of open orders cannot be rebuilt, the venue shall send an
  unsolicited cancel for every order it can still name, and shall say that it
  cannot name the rest.
  \rationale{A member's view converges on the venue's only if the venue reports
  what it has lost. Ceasing to mention an order is not a report of anything.

  This is the fallback beneath \reqref{R-0073}, not an alternative to it: it
  applies where recovery was attempted and could not complete.}
  \uncovered{no scenario makes recovery fail and then checks what the member is sent}
\end{requirement}

\begin{requirement}{R-0021}{A member is not told an order is unknown while the venue is still recovering}
  A request concerning an order shall not be refused as unrecognised before the
  set of open orders has been brought up to date.
  \rationale{Before then the venue does not know whether it holds the order.
  Answering as though it does not is a statement it cannot support, and the
  member acts on it.}
  \uncovered{no scenario asserts that open orders survive a restart}
\end{requirement}

\begin{requirement}{R-0108}{The work needed to bring a component current is bounded}
  A checkpoint shall be taken often enough that the records a recovering
  component must apply after reading one do not exceed a stated bound, and that
  bound shall be a figure the venue can state.
  \rationale{Recovery takes as long as reading the checkpoint plus applying
  everything since it. Leaving the interval between checkpoints unstated leaves
  recovery time unstated with it, and the interval is then set by whatever was
  convenient rather than by how long the venue may be absent.}
  \uncovered{no scenario reads a checkpoint, because none is written}
\end{requirement}

\begin{requirement}{R-0102}{A checkpoint that cannot be trusted is not used}
  A component shall establish that a checkpoint is whole before rebuilding from
  it, and shall not act on the part it could read. A checkpoint that is absent,
  incomplete or inconsistent shall be treated as no checkpoint at all.
  \rationale{A checkpoint is read once, at the moment nothing else knows what the
  component holds, and it is believed entirely. A partial one produces a set of
  orders that is wrong in a way no later check detects, because every subsequent
  check is made against it. Treating a doubtful checkpoint as absent gives the
  answer of \reqref{R-0020}, which is a worse outcome for a member and a knowable
  one.}
  \covers{ha_test:52}
\end{requirement}

\begin{requirement}{R-0022}{The venue keeps enough history to bring a recovered set of orders up to date}
  The record of committed orders shall be retained at least as far back as the
  oldest checkpoint any component may restart from.
  \rationale{Recovery that reads a checkpoint and then asks for what followed it
  fails entirely if what followed has been discarded, and fails in the way that
  is hardest to notice: the checkpoint loads, the catch-up returns nothing, and
  the set looks complete.}
  \uncovered{no scenario exhausts retention, which is anchored to nothing measurable}
\end{requirement}
